\section{Results and Discussion}
\label{sec:results_discussion}

The results demonstrate that the main robot-side functions can be executed on
the STM32 without a host computer. The strength of the evidence differs between
subsystems: computational benchmarks used repeated measurements, the controller
comparison used one hardware run per method, and the E-Skin validation verified
the sensor and UART data path but did not quantify the complete dynamic stop
response. The following discussion therefore distinguishes functional
demonstration from accuracy, repeatability, and safety claims.

\subsection{Embedded Motion Stack}

The STM32 communicated directly with the Feetech bus at
\(1\,\mathrm{Mbit/s}\), and the implemented protocol covered torque commands,
position targets, and encoder feedback. Clearing stale receive data, validating
packet fields, retrying failed reads, and propagating result codes addressed the
communication failures observed during bring-up. The complete communication
chain and the kinematic motion interfaces were exercised on the physical arm.

Forward and inverse kinematics, joint-limit enforcement, and blocking motion
were combined successfully in the four-waypoint square demonstration. This
validates the functional path from a Cartesian target to checked servo commands
and back to encoder-based pose estimation. It does not, however, demonstrate a
geometrically straight TCP trajectory. Each corner is a waypoint, and no path
interpolator constrains the motion between consecutive targets. The initial
approximately 50-tick joint acceptance criterion also characterizes servo-level
arrival rather than physical Cartesian accuracy.

\subsection{Feedback-Control Performance}

The one-shot baseline left model-based terminal TCP errors between
\(4.21\,\mathrm{mm}\) and \(16.44\,\mathrm{mm}\). With Cartesian resolved-rate
feedback, all four reported values lay between \(0.81\,\mathrm{mm}\) and
\(0.93\,\mathrm{mm}\), as shown in
Table~\ref{tab:controller_endpoint_error}. The mean decreased from
\(10.13\,\mathrm{mm}\) to \(0.87\,\mathrm{mm}\), a reduction of approximately
\(91.4\,\%\). This shows the practical benefit of re-reading the joints and
correcting the coupled Cartesian error instead of relying on one converted IK
target.

The \(1\,\mathrm{mm}\) controller stopping criterion is stricter than the
\(5\,\mathrm{mm}\) tolerance used for the standalone IK solver. The two values
serve different purposes: the latter limits numerical solution effort before a
move, whereas the outer loop refines the position reconstructed from measured
encoder ticks. A larger proportional gain and an additional filtered derivative
term produced more corrective motion near the waypoint, so the simpler
speed-limited P controller was retained. These observations are qualitative;
only the final P controller and one-shot baseline have tabulated endpoint data.

The numerical values must also be interpreted as model consistency rather than
physical TCP accuracy. Target and measured positions are evaluated through the
same forward model, so backlash, compliance, calibration error, and tool
displacement remain unobserved. Furthermore, one run per method is insufficient
to establish statistical repeatability or overshoot.

\subsection{Computational Results}

The specialized homogeneous-transform multiplication was the effective
optimization. It reduced an isolated matrix product from 4608 to 1146 cycles,
corresponding to a \(4.02\)-fold speedup without measured numerical deviation.
In the complete algorithms, forward kinematics improved from 39,142 to 18,889
cycles and inverse kinematics from 884,852 to 438,841 cycles, giving speedups of
\(2.07\) and \(2.02\), respectively.

CMSIS-DSP reduced the isolated trigonometric cost but did not materially improve
the full calculation and introduced a small approximation error. Retaining
standard \texttt{sinf()} and \texttt{cosf()} with the specialized matrix product
therefore provided the best measured trade-off. All 100 benchmark IK cases
converged in an average of 5.2 iterations, and the maximum residual error of
\(4.49\,\mathrm{mm}\) remained within the configured solver tolerance. These
benchmarks isolate computation from communication and mechanical delays; they
do not predict complete motion duration.

\subsection{Sensing and System Integration}

The sensing subsystem operated with two physical E-Skin nodes, while the ESP32
software supports up to 32 sensor IDs. Aggregating the latest measurements into
one maximum value removed conflicting LED decisions between nodes. The ESP32
visual feedback uses its own warning, clear, and touch thresholds, whereas the
STM32 independently applies \(0.050\) for motion pause and \(0.020\) for resume
qualification.

The ESP32-to-STM32 UART path was validated on hardware with continuous valid
frames. Unloaded values around \(0.011\) and proximity values above \(0.8\) were
observed. On the STM32, interrupt-driven reception, foreground parsing, startup
validation, stale-data detection, hold commands, and hysteretic resume logic
are integrated into the motion path. Nevertheless, the transmitted maximum is
limited to 4 Hz, and no repeated measurement of detection-to-hold latency during
arm motion is reported. The result is therefore a physically connected and
firmware-integrated collision-feedback prototype, not a validated
collision-avoidance or emergency-stop system.

\subsection{Limitations and Next Steps}

The kinematic model controls position with four active joints; wrist rotation
and the gripper remain fixed, so full six-degree-of-freedom pose control is not
provided. The project also lacks general trajectory planning, guaranteed
straight-line interpolation, systematic workspace and singularity mapping, a
dynamic robot model, and exhaustive unreachable-target testing. Blocking servo
transactions make control timing dependent on communication, while the E-Skin
response remains software-based and leaves the actuators energized.

The most useful next steps are repeated controller trials with independent TCP
metrology, measured sensor-to-hold latency, systematic reachable-workspace and
singularity tests, and explicit Cartesian path generation. Full-pose control
would additionally require incorporating wrist rotation and defining the
desired end-effector orientation. Any safety claim would require independent
hardware measures and a dedicated risk assessment beyond this proof of concept.
